%----------------------------------------------------------------------------------------
%	PACKAGES AND THEMES
%----------------------------------------------------------------------------------------
\documentclass[aspectratio=169,xcolor=dvipsnames]{beamer}
\usetheme{GeminiUOE}

\usepackage{hyperref}
\usepackage{amsmath}
\usepackage{amsfonts}
\usepackage{amssymb}
\usepackage{mathtools}
\usepackage{multicol}
\usepackage{bbm}
\usepackage{amsthm}
\usepackage{tikz}
\usepackage{enumerate}
\usepackage{abbreviations}
%\usepackage{enumitem}
\usepackage{cprotect}
\usepackage{yhmath}
\usepackage{caption}
\usepackage{subcaption}
\usepackage{float}
\usepackage{placeins}
\usepackage{graphicx} % Allows including images
\usepackage{booktabs} % Allows the use of \toprule, \midrule and \bottomrule in tables

%----------------------------------------------------------------------------------------
%	TITLE PAGE
%----------------------------------------------------------------------------------------

\title[short title]{How to Win the Board Game 
\\``Quacks of Quedlinburg''} % The short title appears at the bottom of every slide, the full title is only on the title page
\subtitle{Year 4 Project}

\author[ER, JB, SL] {Eric Rogers, Joseph Brudenell, Sam Lee}

\institute[UoE] % Your institution as it will appear on the bottom of every slide, may be shorthand to save space
{
    The University of Edinburgh % Your institution for the title page
}
\date{Presentation Date TBC, Presentation Room TBC} % Date, can be changed to a custom date


%----------------------------------------------------------------------------------------
%	PRESENTATION SLIDES
%----------------------------------------------------------------------------------------

\begin{document}

\begin{frame}
  % Print the title page as the first slide
  \titlepage
\end{frame}

%\begin{frame}{Overview}
  % Throughout your presentation, if you choose to use \section{} and \subsection{} commands, these will automatically be printed on this slide as an overview of your presentation
%  \tableofcontents
%\end{frame}

%------------------------------------------------
%\section{First Section}
%------------------------------------------------

\begin{frame}{Introduction} % joseph
 % what we're about to do
\end{frame}

\begin{frame}{The Rules} %  sam 

\end{frame}

\begin{frame}{Modelling} % eric
 % assumptions
 % what a policy is, etc
 % opp modelling
 % monte carlo, estimator, python
 % lil bit of mafs
\begin{itemize}
    \item Assumptions\pause
    \item Policies - 
    \[
        \pi:\calS\times\NN^2\to\calA_\pi
    \]
    \pause    
    \vspace*{-0.9cm}
    \begin{gather*}
        \Downarrow\\
         \boldsymbol{\pi} = (\pi_{\mathrm{draw}}, \pi_{\mathrm{bust}}, \pi_{\mathrm{buy}}, \pi_{\mathrm{ruby}}, \pi_{\mathrm{flask}})
    \end{gather*}
    \pause
    \item Estimation - 
    \[
        \widehat{V}_n(\boldsymbol{\pi}\,|\,M)=\frac{1}{n}\sum_{i=1}^n X_i^{(\boldsymbol{\pi}, M)},\qquad X_1^{\boldsymbol{\pi}, M},\ldots,X_n^{\boldsymbol{\pi},M}\,\,\text{iid}
    \]
    \pause
    \[
    M:S_t\to\NN^2,\qquad M(S_t)=\widehat{B}_t^{\mathrm{opps}}
    \]
    \item Monte Carlo Simulation - \texttt{python}
\end{itemize}
\end{frame}

\begin{frame}{Strategy} % joseph

\end{frame}

\begin{frame}{Testing/Results} % sam

\end{frame}


\begin{frame}{Conclusion} % eric
 % what optimal policy in what circumstance - means
 % additional useful findings
 % few bits of further scope
\end{frame}

%------------------------------------------------

\begingroup
\setbeamertemplate{footline}{} % remove footline
\begin{frame}
  \Huge{\centerline{Questions}}
\end{frame}
\endgroup

%----------------------------------------------------------------------------------------

\end{document}